\documentclass[sigconf,review,anonymous]{acmart}

\settopmatter{printacmref=true}
\pagestyle{plain}
\raggedbottom

\usepackage{booktabs}
\usepackage{graphicx}
\usepackage{tabularx}
\usepackage{array}
\usepackage{multirow}
\usepackage{xspace}

\newcolumntype{Y}{>{\raggedright\arraybackslash}X}
\newcommand{\tool}{AUDIT-SID\xspace}
\newcommand{\sid}{SID\xspace}

\title{\tool: Artifact-Level Diagnostics for Semantic-ID Tokenizers in Generative Recommendation}

\author{Anonymous Local Draft}
\affiliation{%
  \institution{Anonymous Institution}
  \country{}
}

\begin{document}

\begin{CCSXML}
<ccs2012>
 <concept>
  <concept_id>10002951.10003317.10003318</concept_id>
  <concept_desc>Information systems~Recommender systems</concept_desc>
  <concept_significance>500</concept_significance>
 </concept>
 <concept>
  <concept_id>10002951.10003317.10003338</concept_id>
  <concept_desc>Information systems~Retrieval models and ranking</concept_desc>
  <concept_significance>300</concept_significance>
 </concept>
</ccs2012>
\end{CCSXML}

\ccsdesc[500]{Information systems~Recommender systems}
\ccsdesc[300]{Information systems~Retrieval models and ranking}
\keywords{semantic IDs, generative recommendation, tokenizer diagnostics, recommender systems}

\begin{abstract}
Semantic-ID tokenizers are a core component of generative recommendation:
items are mapped to discrete code sequences and retrieved by sequence
generation. Yet most evaluations still emphasize final ranking metrics, which
can hide codebook collapse, full-code collisions, collaborative-neighborhood
mismatch, tail compression, or large prefix structures. We present
\tool, a mapping-first diagnostic toolkit for item-to-\sid tokenizer
artifacts. \tool standardizes item mappings, metadata, interactions, and
optional generator outputs, then reports five main artifact diagnostics:
utilization, collision profile, semantic-collaborative alignment, head-tail
capacity allocation, and structural cost proxies. In public-resource
demonstrations, \tool ingests a GRID/RQ-KMeans-style export and a bounded
ReSID/GAOQ export, with method-inspired controllers for collision
qualification, capacity pressure, and variable-depth cost. A
same-item Musical case study on 23,742 items exposes sharply different artifact profiles: a
GRID feature-text residual-k-means row has 3,749 unique full \sid{}s and a
0.9769 full-collision rate, while the bounded ReSID/GAOQ row is
collision-free under its exported mapping; rerunning the GRID row across three
seeds keeps full-collision rate in the narrow 0.9751--0.9769 range. The contribution is a reusable
artifact audit suite, not a new tokenizer or a full generative-recommender
leaderboard.
\end{abstract}

\maketitle

\section{Introduction}

Semantic-ID (\sid) tokenization has become a practical design point for
generative recommendation. Systems such as TIGER map items into discrete code
sequences so that a sequence model can generate candidate identifiers rather
than score every catalog item directly~\cite{rajput2023tiger}. This shift has
made the tokenizer itself a first-order artifact. Recent work explores
residual quantization frameworks, recommender-native encoding, differentiable
\sid learning, non-uniform quantization, qualified collisions, and continual
tokenization~\cite{ju2025grid,liang2026resid,fu2026diger,wei2026card,hu2026quasid,feng2026dact}.
The broader identifier literature also shows that item indexing choices matter
before generative retrieval is trained end-to-end~\cite{hua2023indexids}, that
semantic IDs can improve industrial ranking generalization~\cite{singh2023bettersemanticids},
and that collaborative tokenization methods explicitly optimize behavioral
alignment and assignment diversity~\cite{zhu2024cost,wang2024letter}.

The evaluation practice has not caught up with this artifact boundary. Prior
recommender tooling has argued for behavioral tests beyond NDCG-style
aggregates~\cite{chia2022reclist}, but SID tokenizers introduce a distinct
artifact: the generated item address space. A
tokenizer can support a usable downstream Recall@K while still hiding failure
modes that are hard to diagnose from aggregate ranking metrics: underused
codebooks, repeated full codes, semantically neat but behaviorally poor
prefixes, tail items sharing too little capacity, or prefix structures that
increase serving cost. These issues matter because \sid{}s are not merely
features. They define the address space exposed to the generator.

We present \tool, a resource for auditing item-to-\sid{} tokenizer artifacts
before or alongside downstream recommendation evaluation. The goal is
deliberately narrower than proposing a new tokenizer. \tool asks a concrete
question: given a public tokenizer artifact that maps items to code sequences,
what can be said about its capacity use, collisions, behavioral alignment,
head-tail allocation, and deployment cost without retraining a full generator?

The resource contribution is threefold. First, \tool defines a small
mapping-first interface for \sid assignments, item metadata, interaction
histories, and optional generator outputs. Second, it implements D1--D5a
diagnostics for utilization, collisions, semantic-collaborative alignment,
head-tail capacity, and structural serving-cost proxies, with D6 churn support
for continual-tokenizer settings. Third, it provides public resource-demo
evidence on GRID/RQ-KMeans-style, ReSID/GAOQ, and controlled sanity artifacts,
including a same-item Musical case study. The paper follows the CIKM Resource
four-page constraint, where appendices count toward the limit and detailed
logs are therefore placed in the online artifact~\cite{cikm2026resource}.

\section{Toolkit and Diagnostics}

\tool treats a tokenizer export as an artifact, not as an opaque component of a
trained recommender. The minimum input is a table
\texttt{sid\_assignments(item\_id, sid\_0, ..., sid\_L)}. Optional tables
provide item metadata, interaction histories, and generator outputs. This
contract lets adapters normalize heterogeneous repositories while keeping the
audit independent of a particular training loop.

\paragraph{D1: utilization.}
\tool reports per-level usage, prefix counts, entropy-style imbalance
summaries, and dead or underused code indicators. These expose codebook
collapse and uneven capacity even when full ranking metrics are unavailable.

\paragraph{D2: collision profile.}
\tool measures duplicate full \sid{}s and prefix collisions at each depth. In
the current resource version, D2 is a collision profile rather than a causal
harm estimate. Interaction-qualified collision harm, as motivated by
qualification-aware collision work~\cite{hu2026quasid}, is reserved for the
next artifact version.

\paragraph{D3: semantic-collaborative alignment.}
Semantic grouping and collaborative usefulness need not agree. \tool therefore
computes a co-occurrence-based collaborative neighborhood proxy and asks how
often prefix neighborhoods recover those behaviorally related items. Metadata
purity is retained only as an auxiliary diagnostic, not as the definition of
collaborative quality.

\paragraph{D4: head-tail capacity.}
\tool splits items into popularity buckets and measures whether full \sid
capacity and prefix structure are allocated differently across head, mid, and
tail items. This helps separate genuine tokenizer behavior from simply
restating the catalog popularity distribution.

\paragraph{D5a and D6.}
D5a reports structural cost proxies: \sid length, prefix fan-out, duplicate
codes, and trie-like expansion pressure. These are not real generator serving
latencies. If generator outputs are available, future D5b/D7 metrics can add
invalid paths, next-token entropy, and generated-candidate duplication. D6
computes churn between tokenizer refreshes and is useful for drift-aware
continual-tokenization settings~\cite{feng2026dact}.

\begin{table*}[t]
\caption{Method coverage and artifact role in the current \tool resource.}
\label{tab:method-coverage}
\small
\begin{tabularx}{\textwidth}{@{}llllY@{}}
\toprule
Line & Cluster & Dataset & Artifact & Paper role and caveat \\
\midrule
GRID/RQ-KMeans & A & All\_Beauty; Musical & real export &
Main A evidence; Musical row uses processed feature text, not raw-text TIGER. \\
ReSID/GAOQ & B & Musical & real export &
Main B evidence; bounded one-epoch FAMAE and GAOQ export. \\
Sanity tokenizers & control & Musical; MovieLens smoke & deterministic &
Metric non-redundancy controls; not named methods. \\
CARD compact & B/control & Musical; Sports & proxy/stressor &
Backlog or stressor only; not faithful CARD evidence. \\
DACT & drift & Tools & bundled arrays &
Optional D6 churn demo; not a Cluster B replacement. \\
DIGER/CapsID/AdaSID/AsymRec & future & not run & not claimed &
Coverage context only; verify releases before submission. \\
\bottomrule
\end{tabularx}
\end{table*}

\section{Resource Demo and Case Study}

The demo is designed as a diagnostic-resource check rather than a leaderboard.
It asks whether \tool can ingest real public \sid artifacts from both a
canonical residual-quantization line and a recent recommender-native line, then
produce comparable audit tables with explicit caveats. GRID provides the
open RQ-KMeans-style path~\cite{ju2025grid}; ReSID provides the GAOQ line and
processed Musical artifacts~\cite{liang2026resid}. Sanity tokenizers are kept
as controlled references for metric sensitivity.

\begin{table}[t]
\caption{Same-item Musical diagnostic case study. Lower collision is better;
D3 is a co-occurrence neighborhood proxy.}
\label{tab:musical-diagnostic}
\small
\setlength{\tabcolsep}{3.5pt}
\begin{tabular}{@{}lrrrrr@{}}
\toprule
System & Unique & Dup. & Full coll. & D3 L1 & Tail cap. \\
\midrule
GRID feature-text & 3,749 & 0.8421 & 0.9769 & 0.0552 & 0.3695 \\
ReSID GAOQ & 23,742 & 0.0000 & 0.0000 & 0.1535 & 1.0000 \\
\bottomrule
\end{tabular}
\end{table}

Table~\ref{tab:musical-diagnostic} compares two exports on the same 23,742
Musical items. The GRID row is intentionally labeled as feature-text: it uses
processed feature text from the ReSID artifact path and the official
MiniBatchKMeans-style residual quantization module, not a faithful raw-text
TIGER reproduction. Under that controlled setup, the row exposes high collision
pressure and limited tail capacity. The ReSID/GAOQ row exports one full code
per item in this bounded run, yielding zero full-code collisions and higher
D3-L1 collaborative prefix recovery.

This contrast is useful precisely because the table is not a downstream
ranking claim. It shows that the same audit interface can surface different
artifact mechanisms: compact residual-k-means codes can collapse many items
into shared full \sid{}s; recommender-native GAOQ can avoid full collisions in
the exported mapping; and the head-tail table makes the capacity allocation
visible rather than inferred from item popularity. Additional artifact tables
cover sanity controls, GRID scale checks, DACT churn, and a MovieLens schema
smoke, but those are placed in the repository rather than the four-page PDF.

\section{Availability and Limitations}

The artifact repository contains the \tool Python package, adapter scripts,
metric runners, generated CSV, Markdown, and \LaTeX{} tables, provenance logs,
and
the BibTeX/source audit used for this draft. Large local artifacts are kept out
of git and described through manifests so that reviewers can distinguish
public reproducibility assets from local caches.

The current evidence supports a conservative resource-demo claim. It does not
support several stronger claims. \tool does not reproduce TIGER end-to-end; the
GRID Musical row is a controlled feature-text export. The ReSID Sports balanced
GAOQ path was stopped because constrained k-means was CPU-bound under the
available window, so the main ReSID evidence is the bounded Musical export.
CARD is repaired only as a fallback/stressor path and should not be presented
as faithful CARD reproduction~\cite{wei2026card}. D2 is a collision profile,
not strict causal harm. D3 is a diagnostic proxy and is not yet proven
monotonic with Recall@K or NDCG. D5a is a structural cost proxy without real
generator-output latency, and D6 remains optional drift evidence.

These boundaries are intentional. The paper contributes a reusable audit
surface for tokenizer artifacts, not a new \sid tokenizer, a complete
generative-recommender evaluation, or an industrial online-impact study.
Future versions should add interaction-qualified collision harm, generator
output diagnostics, stronger same-dataset method coverage, and unified search
and recommendation artifact checks~\cite{penha2025jointsid}. Industrial SID
deployment studies further motivate tracking these artifact properties as
first-class engineering objects~\cite{ju2026snapchatsid}.


\clearpage
\bibliographystyle{ACM-Reference-Format}
\bibliography{references}

\section*{GenAI Usage Disclosure}
The authors used generative AI tools for drafting assistance, code-review
support, and experiment-log organization. The reported claims, tables, and
citations were checked against the artifact repository and cited sources during
paper preparation.

\end{document}
